\label{sec:partisan-asymmetry}

\subsection{Partisan Split of Competitive Algorithmic Districts}

Table~\ref{tab:partisan-split} shows the partisan split of competitive
algorithmic districts by lean classification.  Among the 85 competitive
algorithmic districts:
\begin{itemize}
  \item 43 lean Democratic ($D_i \geq 0.50$, margin under 10 points).
  \item 42 lean Republican ($D_i < 0.50$, margin under 10 points).
  \item 17 are true toss-ups ($|D_i - 0.5| < 0.025$).
\end{itemize}

\begin{table}[h]
\centering
\caption{Partisan Split of Competitive Districts}
\label{tab:partisan-split}
\begin{tabular}{lcc}
\toprule
Category & Algorithmic & Enacted \\
\midrule
Lean Democratic (D+1 to D+10) & 43 & 31 \\
Lean Republican (R+1 to R+10) & 42 & 34 \\
True toss-up ($\pm$2.5 pts) & 17 & 12 \\
\midrule
D/R ratio (competitive seats) & 1.024 & 0.912 \\
\bottomrule
\end{tabular}
\end{table}

The near-unity D/R ratio (1.024) among competitive algorithmic districts
contrasts with a 0.912 ratio under enacted maps, indicating that enacted
plans have a modest Republican lean in their competitive-seat allocation.
Neither ratio deviates substantially from symmetry, but algorithmic maps
are measurably more balanced.

\subsection{Engaging the Brunell Argument}

\citet{brunell2008redistricting} argues that competitive districts are
normatively inferior because they guarantee that roughly half the district's
voters are represented by a legislator from the opposing party.  Safe seats,
on this view, maximize constituent satisfaction by ensuring that a majority
of each district's residents are represented by their preferred party.

The argument has force but several limitations that the algorithmic context
exposes.  First, the Brunell argument assumes that ``safe'' seats are allocated
equally across parties.  In practice, enacted maps that eliminate competitive
seats tend to produce safe seats unevenly: geographic efficiency gaps mean
that packing Democratic voters into fewer, more lopsided seats generates
more Republican seats than a symmetric allocation would
imply~\citep{stephanopoulos2014}.  Algorithmic maps, by contrast, produce
safe seats that reflect the underlying geographic distribution of partisans
rather than strategic packing.

Second, the Brunell argument applies equally to the minority voter in a safe
seat.  A minority Republican voter in a safe Democratic urban district is no
better represented than a minority Democratic voter in a safe Republican
exurban district.  Neither the argument for competitiveness nor the argument
for safe seats provides a fully satisfying answer to the representation of
the political minority within a district.

Third, from a dynamic perspective, competitive districts serve as the
mechanism by which national policy preferences translate into legislative
seat changes.  A legislature of only safe seats requires implausibly large
national swings to change party control.  The near-symmetric 43D/42R split
of competitive algorithmic districts indicates that algorithmic maps preserve
the responsiveness function without systematically advantaging either party.

\subsection{Are Competitive Algorithmic Districts Systematically Biased?}

We test for systematic partisan bias in competitive algorithmic districts
by comparing the expected seat share to the vote share in a statewide
environment.  Under the 2020 presidential baseline (Biden 51.3\% nationwide),
a map with no partisan bias would allocate roughly 51.3\% of competitive
seats to Democrats.  Our algorithmic maps allocate 50.6\% of competitive
seats to Democrats (43/85), a difference of 0.7 points---well within sampling
variation.  We conclude that algorithmic competitive districts do not exhibit
systematic partisan bias.
